% !TEX root = Thesis.tex
\documentclass[a4paper,11pt]{mimosis}

\usepackage{metalogo}

%%%%%%%%%%%%%%%%%%%%%%%%%%%%%%%%%%%%%%%%%%%%%%%%%%%%%%%%%%%%%%%%%%%%%%%%
% Some of my favourite personal adjustments
%%%%%%%%%%%%%%%%%%%%%%%%%%%%%%%%%%%%%%%%%%%%%%%%%%%%%%%%%%%%%%%%%%%%%%%%
%
% These are the adjustments that I consider necessary for typesetting
% a nice thesis. However, they are *not* included in the template, as
% I do not want to force you to use them.

% This ensures that I am able to typeset bold font in table while still aligning the numbers
% correctly.

\usepackage{etoolbox}
\usepackage{physics}

% \usepackage{afterpage} %% for page specific geometry


% Make KOMA script happy?
% fore more details see documentatio of KOMA Script options on margins and textarea at
% mirrors.ctan.org/macros/latex/contrib/koma-script/doc/scrguien.pdf
\KOMAoptions{DIV=12}
\KOMAoptions{BCOR=12mm}


\usepackage{multicol}
\setlength{\columnseprule}{0pt}
\setlength\columnsep{10pt}

\usepackage[absolute,overlay]{textpos}


\newcommand{\PhDname}{Peter Pan}


\newcommand{\PhDTitleEN}{How to never grow up}	

\newcommand{\keywordsEN}{keyword 1, kw2, kw3}

\newcommand{\abstractEN}{
  Some abstract
}



\newcommand{\PhDTitleFR}{bleu, bli, blub}

\newcommand{\keywordsFR}{bleu, bli, blub}

\newcommand{\abstractFR}{
  Quelques mots en français.
}	


%% Numéro National de Thèse (donnée par la bibliothèque à la suite du 1er dépôt)/ National Thesis Number (given by the Library after the first deposit)
\newcommand{\NNT}{20XXIPPAXXXX}
%% Full name of Doctoral School : École doctorale de l'Institut Polytechnique de Paris, École doctorale de mathématiques Hadamard
\newcommand{\ecodoctitle}{de l'Institut Polytechnique de Paris}
%% Sigle de l'ED : EDIPP, EDMH / Acronym of the Doctoral School : EDIPP, EDMH
\newcommand{\ecodocacro}{EDIPP}
%% Doctoral School number : 626 (EDIPP), 574 (EDMH)
\newcommand{\ecodocnum}{626}
%% Spécialité de doctorat / Speciality 
\newcommand{\PhDspeciality}{Physique}
%% Établissement de préparation / PhD working place :  l'École polytechnique, l'École nationale supérieure de techniques avancées, l'École nationale de la statistique et de l’administration économique, Télécom ParisTech, Télécom SudParis, l’École des hautes études commerciales de Paris   
\newcommand{\PhDworkingplace}{l'École Polytechnique}
%%  Place of defense
\newcommand{\defenseplace}{Palaiseau}
%% Date de soutenance / Date of defense
\newcommand{\defensedate}{DATE}



\newcommand{\logoEtt}{\includegraphics[height=30mm]{graphics/frontmatter/logo_CPHT.pdf}}
%% Logo of laboratory or establishment.
\newcommand{\vpos}{0.2}
% If needed, modify to align logos vertilcally
\newcommand{\hpos}{13}
%%  If needed, modify to align logos horizontaly
\newcommand{\logoEt}{\includegraphics[height=50mm]{graphics/frontmatter/X_IPParis_blue.pdf}}
%% Institution logo. Filename correspond to institution acronym : ENSAE, ENSTA, TP, TSP, X
\newcommand{\vpostt}{0.33}
%% If needed, modify to align logos vertilcally
\newcommand{\hpostt}{5.5}
%% If needed, modify to align logos horizontaly


%%% JURY

% Choice of the jury's president is made during the defense. Thus, it must be specified only for the second file deposition in ADUM.
% All the jury members listed here must have been present during the defense.

%  Member n°1 (President)
\newcommand{\jurynameA}{Prénom Nom}
\newcommand{\juryadressA}{Statut, Établissement (Unité de recherche)}
\newcommand{\juryroleA}{Président}

%  Member n°2 (Reviewer)
\newcommand{\jurynameB}{Prénom Nom}
\newcommand{\juryadressB}{Statut, Établissement (Unité de recherche)}
\newcommand{\juryroleB}{Rapporteur}

%  Member n°3 (Reviewer)
\newcommand{\jurynameC}{Prénom Nom}
\newcommand{\juryadressC}{Statut, Établissement (Unité de recherche)}
\newcommand{\juryroleC}{Rapporteur}

%  Member n°4 (Examiner)
\newcommand{\jurynameD}{Prénom Nom}
\newcommand{\juryadressD}{Statut, Établissement (Unité de recherche)}
\newcommand{\juryroleD}{Examinateur}

%  Member n°5 (Thesis supervisor)
\newcommand{\jurynameE}{Prénom Nom}
\newcommand{\juryadressE}{Statut, Établissement (Unité de recherche)}
\newcommand{\juryroleE}{Directeur de thèse}

%  Member n°6 (Thesis co-supervisor)
\newcommand{\jurynameF}{Prénom Nom}
\newcommand{\juryadressF}{Statut, Établissement (Unité de recherche)}
\newcommand{\juryroleF}{Invité}

%  Member n°7 (Guest)
\newcommand{\jurynameG}{Prénom Nom}
\newcommand{\juryadressG}{Statut, Établissement (Unité de recherche)}
\newcommand{\juryroleG}{Invité}

%  Member n°8 (Guest)
\newcommand{\jurynameH}{Prénom Nom}
\newcommand{\juryadressH}{Statut, Établissement (Unité de recherche)}
\newcommand{\juryroleH}{Invité}

%  More jury members can be added according to the same model

% DELETE THIS FOR use in actual thesis. This just creates filler text
\usepackage{lipsum}

%%%%%%%%%%%%%%%%%%%%%%%%%%
%% GEOMETRY
%%%%%%%%%%%%%%%%%%%%%%%%%%
\usepackage{geometry}
\geometry{
	inner=25mm,
	top=25mm,
	outer=30mm,
	bottom=25mm
}



%%%%%%%%%%%%%%%%%%%%%%%%%%%%%%%%%%%%%%%%%%%%%%%%%%%%%%%%%%%%%%%%%%%%%%%%
% Hyperlinks & bookmarks
%%%%%%%%%%%%%%%%%%%%%%%%%%%%%%%%%%%%%%%%%%%%%%%%%%%%%%%%%%%%%%%%%%%%%%%%

\usepackage[%
  colorlinks = true,
  citecolor  = BrickRed,
  linkcolor  = RoyalBlue,
  urlcolor   = RoyalBlue,
  unicode,
  pdftitle={\PhDTitleEN},
  pdfauthor={\PhDname},
  pdfsubject={Manuscrit de thèse de doctorat},
  ]{hyperref}

\usepackage{bookmark}
\usepackage[capitalise]{cleveref}

%%%%%%%%%%%%%%%%%%%%%%%%%%%%%%%%%%%%%%%%%%%%%%%%%%%%%%%%%%%%%%%%%%%%%%%%
% Bibliography
%%%%%%%%%%%%%%%%%%%%%%%%%%%%%%%%%%%%%%%%%%%%%%%%%%%%%%%%%%%%%%%%%%%%%%%%
%
% I like the bibliography to be extremely plain, showing only a numeric
% identifier and citing everything in simple brackets. The first names,
% if present, will be initialized. DOIs and URLs will be preserved.

\usepackage[%
  autocite     = plain,
  backend      = biber,
  % doi          = true,
  eprint       = true,
  url          = false,
  giveninits   = true,
  hyperref     = true,
  maxbibnames  = 99,
  maxcitenames = 99,
  sorting      = none,
  % sortcites    = true,
  % sortsets     = true,
  style        = phys,
  % citestyle    = authoryear-comp,
  % bibstyle     = phys,
  backref,%
  natbib       = true,
  ]{biblatex}

%%%%%%%%%%%%%%%%%%%%%%%%%%%%%%%%%%%%%%%%%%%%%%%%%%%%%%%%%%%%%%%%%%%%%%%%
% Some adjustments to make the bibliography more clean
%%%%%%%%%%%%%%%%%%%%%%%%%%%%%%%%%%%%%%%%%%%%%%%%%%%%%%%%%%%%%%%%%%%%%%%%
%
% The subsequent commands do the following:
%  - Removing the month field from the bibliography
%  - Fixing the Oxford commma
%  - Suppress the "in" for journal articles
%  - Remove the parentheses of the year in an article
%  - Delimit volume and issue of an article by a colon ":" instead of
%    a dot ""
%  - Use commas to separate the location of publishers from their name
%  - Remove the abbreviation for technical reports
%  - Display the label of bibliographic entries without brackets in the
%    bibliography
%  - Ensure that DOIs are followed by a non-breakable space
%  - Use hair spaces between initials of authors
%  - Make the font size of citations smaller
%  - Fixing ordinal numbers (1st, 2nd, 3rd, and so) on by using
%    superscripts

% Remove the month field from the bibliography. It does not serve a good
% purpose, I guess. And often, it cannot be used because the journals
% have some crazy issue policies.
\AtEveryBibitem{\clearfield{month}}
\AtEveryCitekey{\clearfield{month}}

% Fixing the Oxford comma. Not sure whether this is the proper solution.
% More information is available under [1] and [2].
%
% [1] http://tex.stackexchange.com/questions/97712/biblatex-apa-style-is-missing-a-comma-in-the-references-why
% [2] http://tex.stackexchange.com/questions/44048/use-et-al-in-biblatex-custom-style
%
\AtBeginBibliography{%
  \renewcommand*{\finalnamedelim}{%
    \ifthenelse{\value{listcount} > 2}{%
      \addcomma
      \addspace
      \bibstring{and}%
    }{%
      \addspace
      \bibstring{and}%
    }
  }
}

% Suppress "in" for journal articles. This is unnecessary in my opinion
% because the journal title is typeset in italics anyway.
\renewbibmacro{in:}{%
  \ifentrytype{article}
  {%
  }%
  % else
  {%
    \printtext{\bibstring{in}\intitlepunct}%
  }%
}

% Remove the parentheses for the year in an article. This removes a lot
% of undesired parentheses in the bibliography, thereby improving the
% readability. Moreover, it makes the look of the bibliography more
% consistent.
\renewbibmacro*{issue+date}{%
  \setunit{\addcomma\space}
    \iffieldundef{issue}
      {\usebibmacro{date}}
      {\printfield{issue}%
       \setunit*{\addspace}%
       \usebibmacro{date}}%
  \newunit}

% Delimit the volume and the number of an article by a colon instead of
% by a dot, which I consider to be more readable.
\renewbibmacro*{volume+number+eid}{%
  \printfield{volume}%
  \setunit*{\addcolon}%
  \printfield{number}%
  \setunit{\addcomma\space}%
  \printfield{eid}%
}

% Do not use a colon for the publisher location. Instead, connect
% publisher, location, and date via commas.
\renewbibmacro*{publisher+location+date}{%
  \printlist{publisher}%
  \setunit*{\addcomma\space}%
  \printlist{location}%
  \setunit*{\addcomma\space}%
  \usebibmacro{date}%
  \newunit%
}

% Ditto for other entry types.
\renewbibmacro*{organization+location+date}{%
  \printlist{location}%
  \setunit*{\addcomma\space}%
  \printlist{organization}%
  \setunit*{\addcomma\space}%
  \usebibmacro{date}%
  \newunit%
}

% Display the label of a bibliographic entry in bare style, without any
% brackets. I like this more than the default.
%
% Note that this is *really* the proper and official way of doing this.
\DeclareFieldFormat{labelnumberwidth}{#1\adddot}

% Ensure that DOIs are followed by a non-breakable space.
\DeclareFieldFormat{doi}{%
  \mkbibacro{DOI}\addcolon\addnbspace
    \ifhyperref
      {\href{http://dx.doi.org/#1}{\nolinkurl{#1}}}
      %
      {\nolinkurl{#1}}
}

% Use proper hair spaces between initials as suggested by Bringhurst and
% others.
\renewcommand*\bibinitdelim {\addnbthinspace}
\renewcommand*\bibnamedelima{\addnbthinspace}
\renewcommand*\bibnamedelimb{\addnbthinspace}
\renewcommand*\bibnamedelimi{\addnbthinspace}

% Make the font size of citations smaller. Depending on your selected
% font, you might not need this.
\renewcommand*{\citesetup}{%
  \biburlsetup
  \small
}

\DeclareLanguageMapping{english}{english-mimosis}

\addbibresource{Thesis.bib}



%%%%%%%%%%%%%%%%%%%%%%%%%%%%%%%%%%%%%%%%%%%%%%%%%%%%%%%%%%%%%%%%%%%%%%%%
% Fonts
%%%%%%%%%%%%%%%%%%%%%%%%%%%%%%%%%%%%%%%%%%%%%%%%%%%%%%%%%%%%%%%%%%%%%%%%

\ifxetexorluatex
  \usepackage{unicode-math}
  \usepackage{fontspec}

  % \setmainfont{Garamond}
  % \setmainfont{EB Garamond}
  % \setmathfont{Garamond Math}


  % STIX Math Font
  % \setmathfont{STIXTwoMath-Regular}[
  % Extension={.otf},
  % Path=./fonts/STIX2/,
  % Scale=1]

  % \setmainfont{STIXTwoText}[
  %   Extension={.otf},
  %   Path=./fonts/STIX2/,
  %   UprightFont={*-Regular},
  %   BoldFont={*-Bold},
  %   ItalicFont={*-Italic},
  %   BoldItalicFont={*-BoldItalic}
  % ]
  % \setsansfont{texgyreheros}[
  %   Extension={.otf},
  %   Path=./fonts/Heros/,
  %   UprightFont={*-regular},
  %   BoldFont={*-bold},
  %   ItalicFont={*-italic},
  %   BoldItalicFont={*-bolditalic}
  % ]


  % Classical LaTeX Fonts "Latin Modern"
  \setmathfont{latinmodern-math}[
    Extension={.otf},
    Path=./fonts/Latin_Modern/,
    Scale=1
  ]

  \setmainfont{lmroman}[
    Extension={.otf},
    Path=./fonts/Latin_Modern/,
    UprightFont={*-regular},
    BoldFont={*-bold},
    ItalicFont={*-italic},
    BoldItalicFont={*-bolditalic},
    FontFace={db}{n}{Font=*demi-regular,Ligatures=CommonOff},
    FontFace={db}{it}{Font=*demi-oblique,Ligatures=CommonOff},
    FontFace={sc}{n}{Font=*caps-regular,Ligatures=CommonOff},
    FontFace={sc}{it}{Font=*caps-oblique,Ligatures=CommonOff},
  ]

  \setmonofont{lmmono}[
    Scale=MatchLowercase,
    Extension={.otf},
    Path=./fonts/Latin_Modern/,
    UprightFont={*-regular},
    % BoldFont={*-bold},
    ItalicFont={*-italic}
  ]

  \setsansfont{lmsans}[
    Extension={.otf},
    Path=./fonts/Latin_Modern/,
    UprightFont={*-regular},
    BoldFont={*-bold},
    ItalicFont={*-oblique},
    BoldItalicFont={*-boldoblique}
  ]

  % \setfontfamily\mycoverthesefont{Libre Franklin}
  % \setfontfamily\mycoverthesefont{IBM Plex Sans}
  % \setfontfamily\mycoverfont{STIXTwoText}


  % Set fonts for sections, chapters, and page headers
  % options are:
  % normal (main) Font
  % font series: regular, bold, SmallCaps,...
  % font shape: normal, italic, oblique,...
  % \setkomafont{sectioning}{\normalfont\fontseries{db}\fontshape{n}}
  \setkomafont{sectioning}{\normalfont\fontseries{db}\fontshape{it}}
  \setkomafont{pagehead}{\normalfont\fontseries{sc}\fontshape{n}}

\else
  % \usepackage[lf]{ebgaramond}
  % \usepackage{stix2}
  \usepackage{lmodern}
  % \usepackage[garamondx,cmbraces]{newtxmath}
  % \usepackage[oldstyle,scale=0.7]{sourcecodepro}
  \singlespacing
  \usepackage{amsmath}
  \usepackage{amssymb}
  \usepackage{mathtools}
\fi


%%%%% Uncomment for Glossary anc acronyms

\newacronym[description={Singular value decomposition}]{SVD}{SVD}{singular value decomposition}
\newacronym{SNF}{SNF}{Smith normal form}
% \newacronym[description={Topological data analysis}]{TDA}{TDA}{topological data analysis}

\newglossaryentry{LaTeX}{%
  name        = {\LaTeX},
  description = {A document preparation system},
  sort        = {LaTeX},
}

\newglossaryentry{Real numbers}{%
  name        = {\(\mathfrak{R}\)},
  description = {The set of real numbers},
  sort        = {Real numbers},
}

\makeindex
\makeglossaries

%%%%%%%%%%%%%%%%%%%%%%%%%%%%%%%%%%%%%%%%%%%%%%%%%%%%%%%%%%%%%%%%%%%%%%%%
% Ordinals
%%%%%%%%%%%%%%%%%%%%%%%%%%%%%%%%%%%%%%%%%%%%%%%%%%%%%%%%%%%%%%%%%%%%%%%%

\makeatletter
\@ifundefined{st}{%
  \newcommand{\st}{\textsuperscript{\textup{st}}\xspace}
}{}
\@ifundefined{rd}{%
  \newcommand{\rd}{\textsuperscript{\textup{rd}}\xspace}
}{}
\@ifundefined{nd}{%
  \newcommand{\nd}{\textsuperscript{\textup{nd}}\xspace}
}{}
\makeatother

\renewcommand{\th}{\textsuperscript{\textup{th}}\xspace}


%%%%%%%%%%%%%%%%%%%%%%%%%%%%%%%%%%%%%%%%%%%%%%%%%%%%%%%%%%%%%%%%%%%%%%%%
% Incipit
%%%%%%%%%%%%%%%%%%%%%%%%%%%%%%%%%%%%%%%%%%%%%%%%%%%%%%%%%%%%%%%%%%%%%%%%

\title{\PhDTitleEN}
% \subtitle{}
\author{\PhDname}

\begin{document}

%%%% Template files for a Thesis at the CPHT.
%%%% Configure the content of the following frontmatter in Sources/Frontmatter/configuration.tex
\frontmatter
  %%%%%%%%%%%%%%%%%%%%%%%%%%%%%%%%%%%%%%%%%%%%%%%%%%%%%%%%%%%%%%%%%%%%%%%%%%%%%%%%%%%%%%%%%%%%%%%%%
%%% Modèle pour la 1ère de couverture des thèses préparées à l'Institut Polytechnique de Paris, basé sur le modèle produit par Guillaume BRIGOT / Template for front cover of thesis made at Institut Polytechnique de Paris, based on the template made by Guillaume BRIGOT
%%% Mis à jour par Aurélien ARNOUX (École polytechnique)/ Updated by Aurélien ARNOUX (École polytechnique)
%%% Nouvelle mise à jour par Jan SCHNEIDER (CPHT, École polytechnique)/ Further updated by Jan SCHNEIDER (CPHT, École polytechnique)
%%% Les instructions concernant chaque donnée à remplir sont données en bloc de commentaire / Rules to fill this file are given in comment blocks
%%% ATTENTION Ces informations doivent tenir sur une seule page une fois compilées / WARNING These informations must contain in no more than one page once compiled
%%%%%%%%%%%%%%%%%%%%%%%%%%%%%%%%%%%%%%%%%%%%%%%%%%%%%%%%%%%%%%%%%%%%%%%%%%%%%%%%%%%%%%%%%%%%%%%%


%%%%%%%%%%%%%%%%%%%%%%%%%%%%%%%%%%%%%%%%%%%%%%%%%%%%%%%%%%%%%%%%%%%%%%%
%%% Mise en page / Page layout
%%% NE RIEN MODIFIER EXCEPTÉ SI BESOIN / DO NOT MODIFY EXCEPT IF NEEDED
%%%%%%%%%%%%%%%%%%%%%%%%%%%%%%%%%%%%%%%%%%%%%%%%%%%%%%%%%%%%%%%%%%%%%%%%

\newgeometry{margin=1cm}



\begingroup\label{cover}
% \mycoverfont
\thispagestyle{empty}

% \color{black} 
\hfill \vfill \tiny \ecodocnum
\begin{textblock}{5}(0,0)
	\textblockcolour{black}
    \centering
    \vspace{8mm}
	\includegraphics[width=5cm]{graphics/frontmatter/ip-paris-white.pdf}
	\vspace{300mm}
\end{textblock}


\begin{textblock}{1}(0.6,3)
	\Large{\rotatebox{90}{\textcolor{white}{\textsf{\textbf{NNT~:~\NNT}}}}}
\end{textblock}

\begin{textblock}{1}(1,7)
	\rotatebox{90}{\textcolor{white}{{
		\fontsize{44pt}{44pt}\selectfont 
		% \mycoverthesefont
		\textsf{\textbf{Thèse de doctorat}}
		}}
	}
\end{textblock}        


\textblockcolour{white}
\begin{textblock}{1}(\hpostt,\vpostt)
	\textblockcolour{white}
	\logoEtt 
\end{textblock}

\begin{textblock}{1}(\hpos,\vpos)
	\textblockcolour{white}
	\logoEt
\end{textblock}

%\vspace{6cm}
%% Texte
\begin{textblock}{10}(5.5,3.5)
	\textblockcolour{white}
	
	% \color{black}
	%\begin{center}  
	\begin{flushright}
		\huge{\PhDTitleEN} \\ \bigskip %% Titre de la thèse 
		\vfill
		% \color{black} %% Couleur noire du reste du texte
		\normalsize {Thèse de doctorat de l'Institut Polytechnique de Paris} \\
		préparée à \PhDworkingplace \\ \bigskip
		\vfill
		École doctorale n\(^{\circ}\)\ecodocnum~\ecodoctitle~(\ecodocacro)  \\
		
		\small{Spécialité de doctorat: \PhDspeciality} \\ \bigskip %% Spécialité 
		\vfill  
		\footnotesize{Thèse présentée et soutenue à \defenseplace, le \defensedate, par} \\ \bigskip
		\vfill
		\Large{\textbf{\textsc{\PhDname}}} %% Nom du docteur
		\vfill
		%\bigskip
	\end{flushright}
	
	%\end{center}
	% \color{black}
	%% Jury
	\begin{flushleft}
		
		\small Composition du Jury~:
	\end{flushleft}
	%% Members of the jury

	\small
	%\begin{center}
	\newcolumntype{L}[1]{>{\raggedright\let\newline \\ \arraybackslash\hspace{0pt}}m{#1}}
	\newcolumntype{R}[1]{>{\raggedleft\let\newline \\ \arraybackslash\hspace{0pt}}lm{#1}}
	
	
    %% Update if members have been added or removed
	\begin{flushleft}\label{jury}
	\begin{tabular}{@{} L{9cm} R{4cm}}
		\jurynameA  \\ \juryadressA & \juryroleA \\[1em]
		\jurynameB  \\ \juryadressB & \juryroleB \\[1em]
		\jurynameC  \\ \juryadressC & \juryroleC \\[1em]
		\jurynameD  \\ \juryadressD & \juryroleD \\[1em]
		\jurynameE  \\ \juryadressE & \juryroleE \\[1em]
		\jurynameF  \\ \juryadressF & \juryroleF \\[1em]
		\jurynameG  \\ \juryadressG & \juryroleG \\[5pt]
		\jurynameH  \\ \juryadressH & \juryroleH \\[5pt]
	\end{tabular} 
	\end{flushleft}   
	%\end{center}
\end{textblock}
\endgroup
\clearpage

\restoregeometry

  \newgeometry{margin=3cm}
\begin{titlepage}
  \vspace*{3cm}
  \makeatletter
  \begin{center}
    \begin{huge}
      \@title
    \end{huge}\\[1.15em]
    %
    \ifdefined \r@subtitle
    \begin{Large}
      \@subtitle
    \end{Large}\\[0.6em]
    \fi
    %
    \emph{by}\\[0.6em]
    \@author
    %
    % \vfill
    % \begin{center}
    %   \includegraphics[height=10cm]{graphics/frontmatter/X_IPParis_blue.pdf}
    % \end{center}
    %
    \vfill
    
    {\large Thèse de doctorat de l'Institut Polytechnique de Paris\\
    préparée au Centre de Physique Théorique\\
    à\\
    \textsc{l'École Polytechnique}}
   
  \end{center}
  \makeatother
\end{titlepage}

\restoregeometry

\newpage
\null
\thispagestyle{empty}
\newpage

  \selectlanguage{british}
% make abstract in english appear in PDF reader list of contents
\phantomsection{}
\addcontentsline{toc}{chapter}{\protect\numberline{}Abstract}


\vspace*{2ex}

\begin{center}
  {\huge \textit{ {\fontseries{md} Abstract}} }
\end{center}

\vspace*{4ex}


% \newcommand{\new}[1]{\textcolor{OliveGreen}{#1}}

\noindent \lipsum

\newpage
\selectlanguage{french}
% make abstract in french appear in PDF reader list of contents
\phantomsection{}
\addcontentsline{toc}{chapter}{\protect\numberline{}Résumé en français}


% \vspace*{2ex}

\begin{center}
  {\huge \textit{\fontseries{md} Résumé} }
\end{center}

\vspace*{3ex}

\noindent \lipsum



\selectlanguage{british}


  \selectlanguage{english}
  \tableofcontents

\mainmatter

%%%%%%%%%%%%% If you want to Partition your thesis, use the \part command like;
  \part[A good part]{%
    A good part\\
    %
    \vspace{1cm}
    %
    \begin{minipage}[l]{\textwidth}
    %
    \textnormal{%
      \normalsize
      %
      \begin{singlespace*}
        \onehalfspacing
        %
        You can also use parts in order to partition your great work
        into larger `chunks'. This involves some manual adjustments in
        terms of the layout, though.
      \end{singlespace*}
    }
    \end{minipage}
  }

  %%%%%%%%%%%%%%%%%%%%%%%%%%%%%%%%%%%%%%%%%%%%%%%%%%%%%%%%%%%%%%%%%%%%%%%%
\chapter{Template}
%%%%%%%%%%%%%%%%%%%%%%%%%%%%%%%%%%%%%%%%%%%%%%%%%%%%%%%%%%%%%%%%%%%%%%%%
% \minitoc
\begin{center}
  \begin{minipage}{0.5\textwidth}
    \begin{small}
      In which the reasons for creating this package are laid bare for the
      whole world to see and we encounter some usage guidelines.
    \end{small}
  \end{minipage}
  \vspace{0.5cm}
\end{center}


\noindent This package contains a minimal, modern template for writing your
thesis. While originally meant to be used for a Ph.\,D.\ thesis, you can
equally well use it for your honour thesis, bachelor thesis, and so
on---some adjustments may be necessary, though.

%%%%%%%%%%%%%%%%%%%%%%%%%%%%%%%%%%%%%%%%%%%%%%%%%%%%%%%%%%%%%%%%%%%%%%%%
\section{Why?}
%%%%%%%%%%%%%%%%%%%%%%%%%%%%%%%%%%%%%%%%%%%%%%%%%%%%%%%%%%%%%%%%%%%%%%%%

I was not satisfied with the available templates for \LaTeX{} and wanted
to heed the style advice given by people such as Robert
Bringhurst~\cite{Bringhurst12} or Edward R.\
Tufte~\cite{Tufte90,Tufte01}. While there \emph{are} some packages out
there that attempt to emulate these styles, I found them to be either
too bloated, too playful, or too constraining. This template attempts to
produce a beautiful look without having to resort to any sort of hacks.
I hope you like it.

%%%%%%%%%%%%%%%%%%%%%%%%%%%%%%%%%%%%%%%%%%%%%%%%%%%%%%%%%%%%%%%%%%%%%%%%
\section{How?}
%%%%%%%%%%%%%%%%%%%%%%%%%%%%%%%%%%%%%%%%%%%%%%%%%%%%%%%%%%%%%%%%%%%%%%%%

The package tries to be easy to use. If you are satisfied with the
default settings, just add
%
\begin{verbatim}
\documentclass{mimosis}
\end{verbatim}
%
at the beginning of your document. This is sufficient to use the class.
It is possible to build your document using either \LaTeX|, \XeLaTeX, or
\LuaLaTeX. I personally prefer one of the latter two because they make
it easier to select proper fonts.

%%%%%%%%%%%%%%%%%%%%%%%%%%%%%%%%%%%%%%%%%%%%%%%%%%%%%%%%%%%%%%%%%%%%%%%%
\section{Making this template \emph{yours}}
%%%%%%%%%%%%%%%%%%%%%%%%%%%%%%%%%%%%%%%%%%%%%%%%%%%%%%%%%%%%%%%%%%%%%%%%

Prior to using this template, the first thing you want to do is probably
a little bit of customisation. You can achieve quick changes in look and
feel by picking your own fonts. With the \verb|fontspec| package loaded
and  \XeLaTeX or \LuaLaTeX as your compiler, this is pretty simple:
%
\begin{verbatim}
\setmainfont{Your main font}
\setsansfont{Your sans-serif font}
\setmonofont{Your monospaced font}
\end{verbatim}
%
Make sure to select nice combinations of that are pleasing to
\emph{your} eyes---this is your document and it should reflect your own
style. Make sure to specify font names as they are provided by your
system. For instance, you might want to use the following combination:
%
\begin{verbatim}
\setmainfont{Libre Baskerville}
\setsansfont[Scale=MatchLowercase]{IBM Plex Sans}
\setmonofont[Scale=MatchLowercase]{IBM Plex Mono}
\end{verbatim}
%
If these fonts exist on your system, your normal text will look
{\fontspec{Libre Baskerville}{a little bit different from the other font used
in this example PDF}}, while your sans-serif font {\fontspec[Scale=MatchLowercase]{IBM Plex Sans}will 
pair nicely with your} {\fontspec[Scale=MatchLowercase]{IBM Plex Mono}{monospaced font}}.


You can also remove the \verb|Scale| directive, but I find that most
fonts pair better if they are adjusted in size a little bit. Experiment
with it until you finds a combination that you enjoy.

%%%%%%%%%%%%%%%%%%%%%%%%%%%%%%%%%%%%%%%%%%%%%%%%%%%%%%%%%%%%%%%%%%%%%%%%
\section{Features}
%%%%%%%%%%%%%%%%%%%%%%%%%%%%%%%%%%%%%%%%%%%%%%%%%%%%%%%%%%%%%%%%%%%%%%%%

%%%%%%%%%%%%%%%%%%%%%%%%%%%%%%%%%%%%%%%%%%%%%%%%%%%%%%%%%%%%%%%%%%%%%%%%
\begin{table}
  \centering
  \begin{tabular}{ll}
    \toprule
    \textbf{Package}      & \textbf{Purpose}\\
    \midrule
      \texttt{amsmath}          & Basic mathematical typography\\
      \texttt{amsthm}           & Basic mathematical environments for proofs etc.\\
      \texttt{booktabs}         & Typographically light rules for tables\\
      \texttt{bookmarks}        & Bookmarks in the resulting PDF\\
      \texttt{dsfont}           & Double-stroke font for mathematical concepts\\
      \texttt{graphicx}         & Graphics\\
      \texttt{hyperref}         & Hyperlinks\\
      \texttt{multirow}         & Permits table content to span multiple rows or columns\\ 
      \texttt{paralist}         & Paragraph~(`in-line') lists and compact enumerations\\
      \texttt{scrlayer-scrpage} & Page headings\\
      \texttt{setspace}         & Line spacing\\
      \texttt{siunitx}          & Proper typesetting of units\\
      \texttt{subcaption} & Proper sub-captions for figures\\
    \bottomrule
  \end{tabular}
  \caption{%
    A list of the most relevant packages required~(and automatically imported) by this template.
  }
  \label{tab:Packages}
\end{table}
%%%%%%%%%%%%%%%%%%%%%%%%%%%%%%%%%%%%%%%%%%%%%%%%%%%%%%%%%%%%%%%%%%%%%%%%

The template automatically imports numerous convenience packages that
aid in your typesetting process. \autoref{tab:Packages} lists the
most important ones. Let's briefly discuss some examples below. Please
refer to the source code for more demonstrations.

%%%%%%%%%%%%%%%%%%%%%%%%%%%%%%%%%%%%%%%%%%%%%%%%%%%%%%%%%%%%%%%%%%%%%%%%
\subsection{Typesetting mathematics}
%%%%%%%%%%%%%%%%%%%%%%%%%%%%%%%%%%%%%%%%%%%%%%%%%%%%%%%%%%%%%%%%%%%%%%%%

This template uses \verb|amsmath| and \verb|amssymb|, which are the
de-facto standard for typesetting mathematics. Use numbered equations
using the \verb|equation| environment.
%
If you want to show multiple equations and align them, use the
\verb|align| environment:
%
\begin{align}
    V &:= \{ 1, 2, \dots \} \\
    E &:= \big\{ \left(u,v\right) \mid \dist\left(p_u, p_v\right) \leq \epsilon \big\}
\end{align}
%
Define new mathematical operators using \verb|\DeclareMathOperator|.
Some operators are already pre-defined by the template, such as the
distance between two objects. Please see the template for some examples. 
%
Moreover, this template contains a correct differential operator. Use \verb|\diff| to typeset the differential of integrals:
%
\begin{equation}
  f(u) := \int_{v \in \domain}\dist(u,v)\diff{v}
\end{equation}
%
You can see that, as a courtesy towards most mathematicians, this
template gives you the possibility to refer to the real numbers~\(\real \) 
and the domain~\(\domain\) of some function. Take a look at the source for
more examples. By the way, the template comes with spacing fixes for the
automated placement of brackets.

%%%%%%%%%%%%%%%%%%%%%%%%%%%%%%%%%%%%%%%%%%%%%%%%%%%%%%%%%%%%%%%%%%%%%%%%
\subsection{Typesetting text}
%%%%%%%%%%%%%%%%%%%%%%%%%%%%%%%%%%%%%%%%%%%%%%%%%%%%%%%%%%%%%%%%%%%%%%%%

Along with the standard environments, this template offers
\verb|paralist| for lists within paragraphs.
%
Here's a quick example: The American constitution speaks, among others, of
%
\begin{inparaenum}[(i)]
  \item life
  \item liberty
  \item the pursuit of happiness.
\end{inparaenum}
%
These should be added in equal measure to your own conduct. To typeset
units correctly, use the \verb|siunitx| package. For example, you might
want to restrict your daily intake of liberty to \SI{750}{\milli\gram}.

Likewise, as a small pet peeve of mine, I offer specific operators for
\emph{ordinals}. Use \verb|\th| to typeset things like July~4\th
correctly. Or, if you are referring to the 2\nd edition of a book,
please use \verb|\nd|. Likewise, if you came in 3\rd in a marathon, use
\verb|\rd|. This is my 1\st rule.

%%%%%%%%%%%%%%%%%%%%%%%%%%%%%%%%%%%%%%%%%%%%%%%%%%%%%%%%%%%%%%%%%%%%%%%%
\section{Changing things}
%%%%%%%%%%%%%%%%%%%%%%%%%%%%%%%%%%%%%%%%%%%%%%%%%%%%%%%%%%%%%%%%%%%%%%%%

Since this class heavily relies on the \verb|scrbook| class, you can use
\emph{their} styling commands in order to change the look of things. For
example, if you want to change the text in sections to \textbf{bold} you
can just use
%
\begin{verbatim}
  \setkomafont{sectioning}{\normalfont\bfseries}
\end{verbatim}
%
at the end of the document preamble---you don't have to modify the class
file for this. Please consult the source code for more information.


% This ensures that the subsequent sections are being included as root
% items in the bookmark structure of your PDF reader.
\bookmarksetup{startatroot}
\backmatter

  %%%%% Glossary %%%%%
  \begingroup
    \let\clearpage\relax
    \glsaddall
    \printglossary[type=\acronymtype]
    \newpage
    \printglossary
  \endgroup

  %%% Bibliography %%%%
  \printindex
  \printbibliography

  %%%%%%% !!!! The contents of this pages must fit in a single A4 page !!!!

\newgeometry{
left=16mm,
top=27mm,
right=16mm,
bottom=27mm
}


\phantomsection{}
\addcontentsline{toc}{chapter}{\protect\numberline{}Back cover}

% 
\thispagestyle{empty}

\fontsize{10}{10}\selectfont
\setlength{\columnseprule}{0pt}
\setlength\columnsep{10pt}

% \begin{textblock*}{61mm}(16mm,3mm)

\vspace*{-16mm}

\noindent\includegraphics[height=24mm]{graphics/frontmatter/ed-ip-paris.pdf}
% \end{textblock*}




\selectlanguage{french}
%%%Titre de la thèse en français / Thesis title in french
\begin{center}
	\fcolorbox{black}{white}{\parbox{0.98\textwidth}{
	\textbf{Titre:} \PhDTitleFR{} 
	\medskip

	%%%Mots clés en français, séprarés par des ;
	\textbf{Mots clés:} \keywordsFR{} 
	\vspace{-3mm}

	%%% Résumé en français / abstract in french
	\begin{multicols}{2}
	\textbf{Résumé:}
  %%%%%%% !!!! The contents of this pages must fit in a single A4 page !!!!
    \lipsum[1-2]
	\end{multicols}
  \vspace{-3.5mm}
	}}
	\end{center}

\vspace*{0mm}


\selectlanguage{british}
%%Titre de la thèse en anglais / Thesis title in english
\begin{center}
  \fcolorbox{black}{white}{\parbox{0.98\textwidth}{
  \textbf{Title:} \PhDTitleEN{}
  
  \medskip
  
  %%% Keywords in english, separated by ;
  \textbf{Keywords:}  \keywordsEN{}
  \vspace{-3mm}
  \begin{multicols}{2}
  \textbf{Abstract:}
    %%%%%%% !!!! The contents of this pages must fit in a single A4 page !!!!
    \lipsum[1-2]
  \end{multicols}
  \vspace{-3.5mm}
  }}
\end{center}



\hfill
\vfill

\noindent
\begin{minipage}{0.5\linewidth}
  \fontsize{11}{12}\selectfont
  \noindent \textbf{Institut Polytechnique de Paris}\\
  \noindent 91120 Palaiseau, France 
\end{minipage}
\hfill
\begin{minipage}{22mm}
\includegraphics[width=\linewidth]{graphics/frontmatter/ip-paris-petit.pdf}
\end{minipage}

\end{document}

\end{document}
